% Teil II: durchgerechneter Weg, Teile a--c.

\subsection*{Gegeben und gesucht}
Lampe an langem Kabel (Masse vernachlaessigbar), Hoehe ueber dem Boden
$\ell_b=\qty{3}{m}$. Gemessen: 5 Perioden in $\qty{63}{s}$, kleine
Amplitude. Gesucht: (a) Bewegungsgleichung, (b) Loesung fuer kleine
Amplituden ($\sin\varphi\approx\varphi$, Ansatz $\varphi=\varphi_0e^{i\omega t}$),
(c) Deckenhoehe.

\subsection*{Ansatz}
Pendel der Laenge $\ell$ (Fadenlaenge), dessen Schwingungsdauer gemessen
wird. Bewegungsgleichung in Tangentialrichtung; Naherung kleiner Winkel;
$\ell$ aus $T=2\pi\sqrt{\ell/g}$.

\subsection*{Rechnung}
\textbf{a) Bewegungsgleichung.} Fuer die Punktmasse am Faden der Laenge
$\ell$ gilt in Tangentialrichtung
\begin{equation}
  m\ell\ddot\varphi=-mg\sin\varphi
  \;\Rightarrow\;
  \boxed{\;\ddot\varphi+\frac{g}{\ell}\sin\varphi=0\;}.
\end{equation}

\textbf{b) Kleine Amplituden.} Mit $\sin\varphi\approx\varphi$:
\begin{equation}
  \ddot\varphi+\frac{g}{\ell}\varphi=0.
\end{equation}
Einsetzen von $\varphi=\varphi_0e^{i\omega t}$ liefert
$\ddot\varphi=-\omega^2\varphi_0e^{i\omega t}$, also
\begin{equation}
  -\omega^2\varphi+\frac{g}{\ell}\varphi=0
  \;\Rightarrow\;\omega=\sqrt{\frac{g}{\ell}},\qquad
  T=\frac{2\pi}{\omega}=2\pi\sqrt{\frac{\ell}{g}}\;.
\end{equation}

\textbf{c) Deckenhoehe.} Fuer 5 Perioden in 63 s ist
$T=\qty{63}{s}/5=\qty{12.6}{s}$. Daraus
\begin{equation}
  \ell=\frac{g T^2}{4\pi^2}
  =\frac{\qty{9.81}{m/s^2}\cdot(\qty{12.6}{s})^2}{4\pi^2}
  \approx\qty{39.5}{m},
\end{equation}
und die Decke liegt um die H\"ohe der Lampe {\"u}ber dem Boden h{\"o}her:
\begin{equation}
  h=\ell+\ell_b\approx\qty{39.5}{m}+\qty{3}{m}=\boxed{\;\qty{42.5}{m}\;}.
\end{equation}

\begin{ergebnis}
$\ddot\varphi+(g/\ell)\sin\varphi=0$; fuer kleine Winkel
$\omega=\sqrt{g/\ell}$, $T=2\pi\sqrt{\ell/g}$; mit $T=\qty{12.6}{s}$ ist
$\ell\approx\qty{39.5}{m}$ und die Deckenhoehe $\approx\qty{42.5}{m}$.
\end{ergebnis}

\begin{kontrolle}
\textbf{Einheiten:} $\ell=gT^2/(4\pi^2)$: [m/s$^2$]$\cdot$s$^2$ = m. \checkmark

\textbf{Querprobe:} Ein ${}\approx40\,\mathrm m$ langes Pendel hat in $g$
gerade eine Periode von etwa 12 bis 13 s -- erwartete Groessenordnung fuer ein
Kirchenseil. Fuer $\ell=\qty{9.8}{m}$ wuerde $T\approx\qty{6.3}{s}$ und 5
Perioden 31 s brauchen, deutlich weniger als gemessen. \checkmark
\end{kontrolle}

\begin{merksatz}
Fuer kleine Amplituden ist das mathematische Pendel ein harmonischer
Oszillator ($T=2\pi\sqrt{\ell/g}$); aus der Periode folgt die Fadenlaenge
und damit die Deckenhoehe.
\end{merksatz}
